\documentclass[10pt,a4paper]{amsart}
\usepackage[margin=19mm]{geometry}
\usepackage{amsmath,amssymb,mathtools}
\usepackage[scr=rsfs,cal=euler]{mathalfa}
\usepackage{xfrac}
\usepackage[unicode,hidelinks,bookmarks=false]{hyperref}
\newcommand{\ie}{i.e.\ }
\newcommand{\cf}{cf.\ }
\newcommand{\resp}{respectively\ }
\newcommand{\Subsection}{\subsection}
\providecommand{\nobreakdash}{\nobreak\hbox{-}}
\setlength{\emergencystretch}{2em}
\multlinegap0pt
\usepackage{natbib}
\usepackage{caption}
\usepackage{algorithm}\usepackage{algpseudocode}
\usepackage{amssymb}
\usepackage{xcolor}
\usepackage{booktabs}
\usepackage{multirow}
\newtheorem{proposition}{Proposition}
\title{Learning from Imperfect Data: Robust Inference of Dynamic Systems using Simulation-based Generative Model}
\author{Hyunwoo Cho, Hyeontae Jo, Hyung Ju Hwang}
\date{Source: arXiv:2507.10884v2 (CC BY 4.0)}
\begin{document}
\maketitle

\noindent\emph{Source and licence.} Learning from Imperfect Data: Robust Inference of Dynamic Systems using Simulation-based Generative Model. Version arXiv:2507.10884v2 (CC BY 4.0), distributed by arXiv under the Creative Commons Attribution 4.0 International licence, http://creativecommons.org/licenses/by/4.0/, as recorded in the arXiv licence metadata for this exact version and shown on the abstract page https://arxiv.org/abs/2507.10884v2. Full licence notice: \texttt{licenses/arxiv-2507.10884-CC-BY-4.0.txt}.\par\smallskip

\noindent\textit{Editorial scope.} Counted unit: the article's proposition gan2 (source0.tex lines 514--520). The article's complete original proof of it is delivered with it (source0.tex lines 521--556, one \texttt{proof} environment of 36 lines). No mathematics is written, repaired or re-derived: the statement and the proof are the article's own lines, in the article's own notation, and no retained line is substituted or re-flowed.  No further source-local context is needed: every cross-reference of every retained line resolves inside the two delivered spans.  Nothing was omitted from any retained span, so the package carries no omission record.  No retained line contains a citation, so the wrapper carries no bibliography and no bibliography entry.  No retained line contains a figure environment, an image-inclusion command or drawing code, so no image is delivered and the package declares no asset dependency.  Editorial formatting only, in addition: an AMS \texttt{amsart} preamble that restates the article's own declarations and macros used by the retained lines, namely the environments \texttt{proposition} on separate counters, together with the standard packages those lines need.  The retained environments print in this document as Proposition 1 in order of appearance, whereas the article prints the same items with the numbering of its own full document.  The complete native source of the article as distributed in the arXiv e-print for this exact version is bundled unchanged as \texttt{sources/181590/source0.tex}.
\begin{proposition}\label{gan2} Suppose that $\pi_{e}^{G},\pi_{e}^{o}$ satisfy \[ \mathbb{E}_{\mathbf{e}_{n}^{G} \sim \pi_{e}^{G}}[\{e_{n,i}^{G}(t_{j})\}_{i\in I, j \in J}] = \mathbb{E}_{\mathbf{e}_{n}^{o} \sim \pi_{e}^{o}}[\{e_{n,i}^{o}(t_{j})\}_{i\in I, j \in J}] = \mathbf{0}.\] Then, if $\pi_{p}^{G}$, $\pi_{e}^{G}$ are the desired probability density functions (pdfs) that make Wasserstein distance $d(Y^{o},Y^{G})$ equal to zero, i.e.,
\begin{align*}
\sup_{\phi \in \mathrm{Lip}_{1}}\mathbb{E}_{\mathbf{p}_{n}^{G} \sim \pi_{p}^{G}, \mathbf{e}_{n}^{G} \sim \pi_{e}^{G}}\bigg[\phi (\{y_{i}(t_{j};\mathbf{p}_{n}^{G})&+e_{n,i}^{G}(t_{j})\}_{i\in I,j\in J})\bigg]\\
& - \mathbb{E}_{\mathbf{e}_{n}^{o} \sim \pi_{e}^{o}}\bigg[\phi (\{y_{i}(t_{j};\mathbf{p}_{true})+e_{n,i}^{o}(t_{j})\}_{i\in I,j\in J}) \bigg] = 0,
\end{align*}
 then \[\mathbb{E}_{\mathbf{p}_{n}^{G} \sim \pi_{p}^{G}}[y_{i}(t_{j};\mathbf{p}_{n}^{G})] = y_{i}(t_{j};\mathbf{p}_{true})\] and \[ \mathbb{E}_{\mathbf{p}_{n}^{G} \sim \pi_{p}^{G}}[(y_{i}(t_{j};\mathbf{p}_{n}^{G})-y_{i}(t_{j};\mathbf{p}_{true}))^{2}] = \text{Var}(e_{n,i}^{o}(t_{j}))-\text{Var}(e_{n,i}^{G}(t_{j}))\] for all $i \in I$, $j \in J$.
\end{proposition}

\begin{proof}We have that
\begin{align*}
0 &= \int\int \phi(\{y_{i}(t_{j};\mathbf{p}_{n}^{G})+e_{n,i}^{G}(t_{j})\})\pi_{p}^{G}(\mathbf{p}_{n}^{G})
\pi_{e}^{G}(\mathbf{e}_{n}^{G})d\mathbf{p}_{n}^{G}d\mathbf{e}_{n}^{G}\\
 & \quad\quad -\int \phi(\{y_{i}(t_{j};\mathbf{p}_{true})+e_{n,i}^{o}(t_{j})\})\pi_{e}^{o}(\mathbf{e}_{n}^{o})d\mathbf{e}_{n}^{o} \\ 
&= \int\int \phi(x)\pi_{p}^{G}(\mathbf{p}_{n}^{G})\pi_{e}^{G}(\{x_{i,j}-y_{i}(t_{j};\mathbf{p}_{n}^{G})\})d\mathbf{p}_{n}^{G}dx-
\int \phi(x)\pi_{e}^{o}(\{x_{i,j}-y_{i}(t_{j};\mathbf{p}_{true})\})dx \\ &=\int \phi(x)\Bigg(
\int \pi_{e}^{G}(\{x_{i,j}-y_{i}(t_{j};\mathbf{p}_{n}^{G})\})\pi_{p}^{G}(\mathbf{p}_{n}^{G})d\mathbf{p}_{n}^{G}-\pi_{e}^{o}(\{x_{i,j}-y_{i}(t_{j};\mathbf{p}_{true})\})\Bigg)dx \end{align*} for all $\phi \in \text{Lip}_{1}$
and thus we obtain 
\[\mathbb{E}_{\mathbf{p}_{n}^{G} \sim \pi_{p}^{G}}[\pi_{e}^{G}(\{x_{i,j}-y_{i}(t_{j};\mathbf{p}_{n}^{G})\})]=\pi_{e}^{o}(\{x_{i,j}-y_{i}(t_{j};\mathbf{p}_{true})\})\] for any $x_{i,j} \in \mathbb{R}$. \\
Now denote that $\pi_{i,j}^{G}$, $\pi_{i,j}^{o}$ be marginal pdfs for $e_{n,i}^{G}(t_{j})$, $e_{n,i}^{o}(t_{j})$ respectively.

From the given condition, we have that 
\begin{align*}
    \int x_{i,j}\pi_{e}^{G}(\{x_{i',j'}-y_{i'}(t_{j'};\mathbf{p}_{n}^{G})\})dx = y_{i}(t_{j};\mathbf{p}_{n}^{G}),\\ 
\int x_{i,j}\pi_{e}^{o}(\{x_{i',j'}-y_{i'}(t_{j'};\mathbf{p}_{true})\})dx = y_{i}(t_{j};\mathbf{p}_{true}).
\end{align*}
Therefore,
\begin{align*}
\mathbb{E}_{\mathbf{p}_{n}^{G} \sim \pi_{p}^{G}}[y_{i}(t_{j};\mathbf{p}_{n}^{G})] 
&= \mathbb{E}_{\mathbf{p}_{n}^{G} \sim \pi_{p}^{G}}[\int x_{i,j}\pi_{e}^{G}(\{x_{i',j'}-y_{i'}(t_{j'};\mathbf{p}_{n}^{G})\})dx] \\
&= \int x_{i,j}\pi_{e}^{o}(\{x_{i',j'}-y_{i'}(t_{j'};\mathbf{p}_{true})\})dx = y_{i}(t_{j};\mathbf{p}_{true})
\end{align*} 

and 

\begin{align*}
\mathbb{E}_{\mathbf{p}_{n}^{G} \sim \pi_{p}^{G}}&[(y_{i}(t_{j};\mathbf{p}_{n}^{G})-y_{i}(t_{j};\mathbf{p}_{true}))^{2}]\\
&= \mathbb{E}_{\mathbf{p}_{n}^{G} \sim \pi_{p}^{G}}[y_{i}(t_{j};\mathbf{p}_{n}^{G})^{2}]-y_{i}(t_{j};\mathbf{p}_{true})^{2} \\
&= \mathbb{E}_{\mathbf{p}_{n}^{G} \sim \pi_{p}^{G}}[\int x_{i,j}^{2}\pi_{e}^{G}(\{x_{i',j'}-y_{i'}(t_{j'};\mathbf{p}_{n}^{G})\})dx - \text{Var}(e_{n,i}^{G}(t_{j}))]-y_{i}(t_{j};\mathbf{p}_{true})^{2} \\ 
&= (\int x_{i,j}^{2}\pi_{e}^{o}(\{x_{i',j'}-y_{i'}(t_{j'};\mathbf{p}_{true})\})dx - y_{i}(t_{j};\mathbf{p}_{true})^{2}) - \text{Var}(e_{n,i}^{G}(t_{j})) \\ 
&= \text{Var}(e_{n,i}^{o}(t_{j})) - \text{Var}(e_{n,i}^{G}(t_{j}))
\end{align*}

for each $i,j$.
\end{proof}

\end{document}
